\documentclass[../../../main.tex]{subfiles}
\begin{document}

\subsection{Shannon Entropy}
We define the unexpectedness associated with that event as $S(p)$, so that it increases as the event becomes less likely to occur, more unexpectedness in other word.
Thus, the surprise is a decreasing function of the probability.
This quantity is called Shannon information, or Shannon entropy
\begin{equation*}
  S(\left\{ p_i \right\} ) = -\sum_{i=1}^{N} p_i \log p_i
\end{equation*}
This entropy is defined in such way to satisfy
\begin{enumerate}
  \item \textbf{Continuity.}
        It follows from continuity of the logarithm
  \item \textbf{Monotonicity.}
        Less probable events carry larger information, so $S(p)$ is decreasing in $p$
  \item \textbf{Additivity.}
        If two independent events occur with probabilities \( p \) and \( q \) \((0 \leq p, q \leq 1)\), then the surprise when both occur should be the sum of the individual surprises
        \begin{equation*}
          S(p) + S(q) = S(pq)
        \end{equation*}
        It is a functional equation whose continuous solutions are exactly logarithms.
  \item \textbf{Averaging.}
        We impose that the total information is the expectation value of the single-event information
        \begin{equation*}
          S(\left\{ p_i \right\} ) = \sum_{i } p_i  s(p_i)
        \end{equation*}
        Here, the effect of Changing the logarithm base produces only a constant rescaling
        \begin{equation*}
          S_b\bigl(\{p_i\}\bigr) = \frac{1}{\log_a b}\; S_a\bigl(\{p_i\}\bigr)
        \end{equation*}
\end{enumerate}

The definition can be easily confirmed by
\begin{align*}
  S(\{p_i\}, \{q_j\}) & = -\sum_{i,j} p_i q_j \log (p_i q_j) = -\sum_{i,j} p_i q_j \log p_i - \sum_{i,j} p_i q_j \log q_j \\
  S(\{p_i\}, \{q_j\}) & = -\sum_i p_i \log p_i - \sum_j q_j \log q_j = S(\{p_i\}) + S(\{q_j\})
\end{align*}

Let us express this in another way.
If the uncertainty before measurement is large, then the information obtained after measurement, the Shannon information, is also large.
Thus, Shannon information quantifies the uncertainty of the probability distribution before the event occurs.

The asymptotic characterization of Shannon entropy
\begin{equation*}
  H(p) = \frac{1 }{N }\log_2 W
\end{equation*}
represents the degree of incompressibility; which means the smaller the value, the greater the degree of compression.
Here $W= {}_N C_m = 2^{NH(p)}$ refer to the possiblities associated with probability $p$.
The Shannon entropy or $H(p)$ can take the value $0\leq H(p) \leq 1$.
\begin{enumerate}
  \item $H(p)=1$.
        Corresponds to the case of where no compression is possible, meaning all sequences equally likely.
        Consider the case of $N$ bits where all bits have the same probability.
        The entropy in this case is $H[p(N)]=1$, so the number of possible string is $W=2^N$.
        This is the upper bound for representing $N$ bit, meaning no compression is possible.
  \item $0<H(p)<1$.
        Corresponds to the case of biased distribution where $p \neq 1/2$.
        \begin{equation*}
          W \approx 2^{N H(p)} < 2^{N}
        \end{equation*}
        Consequently, the sequence can be encoded using fewer than $N$ bits, specifically $NH(p)$ bits.
        \begin{equation*}
          \log_2 W \approx N H(p) < N
        \end{equation*}
  \item $H(p)=0$.
        Outcome is deterministic, where $p=0$ or $p=1$.
        In this case, $W=1$ and $\log_2 W=0$.
\end{enumerate}

\subsubsection{Example: Coin toss.}
Suppose a coin is tossed, and the probability of obtaining heads is \( p \) (denoted as 0), while the probability of tails is \( 1-p \) (denoted as 1). In this case, the Shannon information \( S(p) \) is
\begin{equation*}
  S(p)
  = -p \log_2 p - (1 - p)\log_2 (1 - p)
  = H(p)
\end{equation*}
The base of the logarithm has been taken as 2 for later convenience; in this case, it is often written as \( H(p) \).

In particular, for \( p = 1/2 \), the Shannon information \( H(p) \) attains its maximum value. When \( p = 0 \) or \( p = 1 \), the Shannon information takes its minimum value of 0.
This corresponds to the situation where the result is completely determined in advance, meaning that no new information is obtained.

\begin{figure}
  \centering
  \begin{tikzpicture}
    \begin{axis}[
        xmin=0, xmax=1,
        ymin=0, ymax=1.1,
        xlabel={$p$},
        ylabel={$H(p)$},
      ]
      \addplot[
        domain=0.001:0.999,
        samples=1000,
        smooth,
      ] {-(x*ln(x) + (1-x)*ln(1-x))/ln(2)};
    \end{axis}
  \end{tikzpicture}
  \caption*{Figure: $H(p)$ with respect to $p$.}
\end{figure}

\subsubsection{Example: Partition function.}
Let the states with energies $\epsilon_i$ $(i = 1,2,\dots,N)$ be realized with probabilities $p_i$.
Under the condition that the average energy is fixed and unitary probability
\begin{equation*}
  \sum_{i=1}^{N} p_i \epsilon_i = E
  \qquad
  \sum_{i=1}^{N} p_i = 1
\end{equation*}
Introduce Lagrange multipliers $\alpha$ and $\beta$, and consider maximizing under this constraints.
Let $p_i$ $(i=1,\dots,N)$ be independent variables, then
\begin{equation*}
  L = -\sum_{i=1}^{N} p_i \log p_i
  + \alpha \left(\sum_{i=1}^{N} p_i - 1\right)
  + \beta \left(E - \sum_{i=1}^{N} p_i \epsilon_i \right)
\end{equation*}
Differentiate with respect to $p_i$ and set it to zero
\begin{equation*}
  \frac{\partial L}{\partial p_i} = -\log p_i - 1 + \alpha - \beta \epsilon_i = 0
\end{equation*}
From this
\begin{equation*}
  p_i = \frac{e^{-\beta \epsilon_i}}{e^{1-\alpha}}= \frac{e^{-\beta \epsilon_i}}{Z}
\end{equation*}
Furthermore, from the constraint condition
\begin{equation*}
  \sum_{i=1}^{N} \frac{e^{-\beta \epsilon_i}}{Z}  =  1\implies
  Z  =  \sum_{i=1}^{N} e^{-\beta \epsilon_i}
\end{equation*}

\subsubsection{Derivation.}
Previously, an intuitive introduction of Shannon information was given.
But, here a derivation closer to Shannon’s original paper.

Messages are generally transmitted in binary form using 0 and 1. For example,
Consider a string of $N$ bits such as $010010110100\ldots$.
Suppose that among the $N$ characters, there are $m$ zeros.
Then, the number of distinct binary sequences that contain exactly $m$ zeros, or $N-m$ ones, is
\begin{equation*}
  W = {}_N C_m = \binom{N}{m} = \frac{N!}{m!(N-m)!}
\end{equation*}
When $N$ and $m$ are large, we use Stirling’s formula
\begin{align*}
  \log W & = N \log N - N - m \log m + m - (N - m)\log (N - m) + N - m                                                      \\
         & = N \log N - m \log m - N\left(1 - \frac{m}{N}\right)\log \left(1 - \frac{m}{N}\right) - (N - m)\log N           \\
  \log W
         & = N \left[ -\frac{m}{N} \log \frac{m}{N} - \left(1 - \frac{m}{N}\right)\log \left(1 - \frac{m}{N}\right) \right] \\
  \ln W
         & =N H(p) \log 2                                                                                                   \\
  W      & = [\exp (\ln 2) ]^{N H(p) } = 2 ^{N H(p) }
\end{align*}
with $p(m)=m/N$ as the probability that 0 appears.

Shannon entropy appears in the binary case as
\begin{equation*}
  H(p)
  =  -\frac{m}{N} \log \frac{m}{N} - \left(1 - \frac{m}{N}\right)\log \left(1 - \frac{m}{N}\right)
\end{equation*}
The step from sequences to probabilities in general $N$-ary sources generalizes this same reasoning to
\begin{equation*}
  S(\left\{ p_i \right\} ) = -\sum_{i=1}^{N} p_i \log p_i
\end{equation*}

\subsection{Shannon Compression}
Shannon coding theorem state that the average number of bits $\braket{l }$ is bounded by $H \le \bar{\ell} < H + 1$.
Shannon proposed the strategy: “represent frequently occurring symbols with a small number of bits, and rarely appearing symbols with a large number of bits, thereby reducing the total number of bits.”
Although in general, this does not achieve the lower bound $H$.
More precisely, any encoding with lengths $l_i$
\begin{equation*}
  \sum_i p_i l_i \ge H(p)
\end{equation*}
with equality if and only if $p_i = 2^{l_i}$.

The algorithm is as follows.
\begin{enumerate}
  \item \textbf{Ordering}.
        Arrange the symbols $a_1,a_2,\ldots,a_i,\ldots,a_n$ in order of decreasing frequency, and let their probabilities be $p_1>p_2>\cdots>p_i>\cdots>p_n$.
  \item \textbf{Probabilities}.
        Compute the $i$-th sum of probabilities from $1$ to $i-1$
        \begin{equation*}
          q_i = \sum_{k=1}^{i-1} p_k
        \end{equation*}
  \item \textbf{Self-information}.
        Compute the surprise (self-information) and take the integer, which is represented using the floor function
        \begin{equation*}
          s_i = \lfloor -\log_2 p_i \rfloor
        \end{equation*}
  \item \textbf{Binary representation.}
        Express $q_i$ in binary fraction form, and take the digits after the decimal point up to the $s_i$-th place.
        This becomes the code for each symbol.
\end{enumerate}

\subsubsection{Example.}
Consider a message consisting of the characters A, B, C, and D, which appear with probabilities 1/2,1/4,1/8,1/8, respectively.
Frequently appearing characters like A are represented with a few bits, and rarely appearing characters like C and D are represented with a larger number of bits, such as
\begin{equation*}
  A = [0] \qquad B = [10]\qquad C = [110]\qquad  D = [111]
\end{equation*}
Compression using this encoding can be done as
\begin{align*}
  [0100101110011010010]
   & \rightarrow [0,10,0,10,111,0,0,110,10,0,10] \\
  [0100101110011010010]
   & \rightarrow \textit{ABABDAACBAB}
\end{align*}
The average bits used is
\begin{equation*}
  \braket{l } = \sum_{i } p_i l_i =
  \frac{1}{2}\cdot 1 + \frac{1}{4}\cdot 2 + \frac{1}{8}\cdot 3 + \frac{1}{8}\cdot 3 = \frac{7}{4} = 1.75
\end{equation*}
The Shannon entropy in this case
\begin{align*}
  H & = 2 -\frac{1}{2}\log_2\left(\frac{1}{2}\right) - \frac{1}{4}\log_2\left(\frac{1}{4}\right) - \frac{1}{8}\log_2\left(\frac{1}{8}\right) - \frac{1}{8}\log_2\left(\frac{1}{8}\right) \\
  H & =  \frac{7}{4} = 1.75
\end{align*}
This is less that the average bits using simple 2-bit encoding, such as
\begin{equation*}
  A = [00] \qquad B = [01] \qquad C = [10]\qquad  D = [11]
\end{equation*}
which obviously have the average of 2.
In this case, $H=1.75$ represent the theoretical minimum average bit length.

\subsubsection{Another example.}
With the given symbols and their probabilities, the average number of bits is
\begin{align*}
  \braket{l }
   & = \frac{1}{2}\cdot 1+\frac{1}{4}\cdot 2+\frac{1}{8}\cdot 3+\frac{1}{16}\cdot 4+\frac{1}{32}\cdot 5+\frac{1}{64}\cdot 6+\frac{1}{128}\cdot 7+\frac{1}{128}\cdot 7 \\
  \braket{l }
   & =  \frac{127}{64} \approx 1.98
\end{align*}
and the Shannon entropy is
\begin{align*}
  H(p) & = -\frac{1}{2}\log_2\left(\frac{1}{2}\right)-\frac{1}{4}\log_2\left(\frac{1}{4}\right)-\frac{1}{8}\log_2\left(\frac{1}{8}\right)-\frac{1}{16}\log_2\left(\frac{1}{16}\right)                \\
       & \qquad -\frac{1}{32}\log_2\left(\frac{1}{32}\right)-\frac{1}{64}\log_2\left(\frac{1}{64}\right)-\frac{1}{128}\log_2\left(\frac{1}{128}\right)-\frac{1}{128}\log_2\left(\frac{1}{128}\right) \\
  H(p) & = \frac{127}{64}\approx 1.98.
\end{align*}
Thus, the average number of bits and the Shannon information coincide.

\begin{table}[]
  \centering
  \caption{Table: Encoding of symbols}
  \begin{tabular}{@{}llllll@{}}
    \toprule
    Symbol & $p_i$ & $q_i$   & $s_i$ & Binary of $q_i$ & Code      \\
    \midrule
    $a_1$  & 1/2   & 0       & 1     & 0               & [0]       \\
    $a_2$  & 1/4   & 1/2     & 2     & 0.1             & [10]      \\
    $a_3 $ & 1/8   & 3/4     & 3     & 0.11            & [110]     \\
    $a_4$  & 1/16  & 7/8     & 4     & 0.111           & [1110]    \\
    $a_5 $ & 1/32  & 15/16   & 5     & 0.1111          & [11110]   \\
    $a_6$  & 1/64  & 31/32   & 6     & 0.11111         & [111110]  \\
    $a_7$  & 1/128 & 63/64   & 7     & 0.111111        & [1111110] \\
    $a_8$  & 1/128 & 127/128 & 7     & 0.1111111       & [1111111] \\
    \bottomrule
  \end{tabular}
\end{table}

\subsection{Law of Large Number}
Roughly speaking, if one takes sufficiently many samples, the average value of the measured physical quantity approaches the expected value of that quantity.
Suppose a physical quantity $A$ is measured independently $N$ times, yielding values $a_1,\dots a_N$.
Also suppose that $P(A)$ is the probability distribution of $A$.
Then for sufficiently large number of measurements $N$, the average $S_N$ approaches the expectation value $E(A)$
\begin{equation*}
S_N=\sum_{i=1}^N\frac{a_i}{N}\rightarrow
E(A)=\int dP(A)\;A 
\end{equation*}
for finite $E(A)$.

\subsubsection{Proof.}
First compute the expectation value of average
\begin{equation*}
E\big(S_N^2\big) = E \left( \frac{1}{N^2}\sum_{i,j} a_i a_j \right) = \frac{1}{N^2}\sum_{i,j} E(a_i a_j)
\end{equation*}
Under the independent and identically distributed (IID) assumption, $E(a_ia_j)=E(a_i)E(a_j)$ for distinct $(i,j)$.
Then
\begin{equation*}
E\big(S_N^2\big) 
= 
\frac{1 }{N^2}\sum_i E(a_i^2) + \sum_{i\neq j} E(a_i a_j)
=\frac{1 }{N^2} \sum_i E(a_i^2) = 
\end{equation*}
where we have assumed $E(a_i)=0$.
Since the variables are identically distributed, $E(a_i^2)=E(A^2)$
\begin{equation*}
E\big(S_N^2\big) 
= 
\frac{1 }{N^2} N E(A^2)
=
\frac{1 }{N} E(A^2)
\end{equation*}

Using the probability distribution, we can also evaluate
\begin{equation*}
  E(S_N^2) =\int dP\;S_N^2=\int_{|S_N|>\epsilon} dP\,S_N^2+\int_{|S_N|\le\epsilon} dP\;S_N^2
\end{equation*}
By definition of $|S_N| > \epsilon$ region, we have $S_N^2 > \epsilon^2$.
We can rewrite second term as 
\begin{equation*}
  \int_{|S_N|>\epsilon} dP\;S_N^{2} \ge \epsilon ^{2}\int_{|S_N|>\epsilon} \epsilon^{2}dP = \epsilon ^2 p(|S_N|>\epsilon)
\end{equation*}
Therefore
\begin{equation*}
 p(|S_N|>\epsilon) < \frac{E(A^2)}{\epsilon^2 N}
\end{equation*}
Since the right-hand side of this expression goes to zero in the limit $N \rightarrow \infty$ with $\epsilon$ fixed, the probability that the sample mean deviates from zero (which is the expectation in this special case) by more than $\epsilon$.
This completes the proof in the case where the expectation value $E(A)$ is zero.
For non-zero case, consider $A-E(A)$ in place of $A$ and, as $N \rightarrow \infty$, one obtain
\begin{equation*}
 p(|S_N-E(A)|>\epsilon) =0
\end{equation*}
This means that the probabilities of $S_N$ deviate by $\epsilon$ from $E(A)$ approaches zero.

\subsection{Type of Entropy}
\subsubsection{Relative entropy.}
Suppose one had been told in advance that the probability distribution is \( \{q_i\} \), but in reality it is \( \{p_i\} \). 
Thus, the difference between the two is \( \log_2 (p_i/q_i) \). 
The average of this with respect to the actual probability distribution \( \{p_i\} \) is
\begin{equation*}
  S(p||q) = \sum_i p_i \log_2 \frac{p_i}{q_i}
\end{equation*}
This is called the relative entropy, which represents the average difference between the two probability distributions $\{p_i\}, \{q_i\} $. 
\subsubsection{Joint entropy.}
First, let $i \in I$ and $j \in J$ be events, and denote by $p(i,j)$ the joint probability that both occur, that is $I$ takes value $i$ and $J$ takes value $j$. 

Then the joint entropy is given by
\begin{equation*}
  S(I, J) = - \sum_{ij} p(i,j) \log_2 p(i,j)
\end{equation*}
This represents the uncertainty before observation of the distribution obtained when $(i,j)$ is regarded as a single combined event.

\subsubsection{Conditional entropy.}
This can be interpreted as the average surprise after knowing $J$. 
In other words, it is the uncertainty under the condition $J$.
\begin{equation*}
S(I|J) = S(I, J) - S(J)
\end{equation*}
hat is the remaining uncertainty, in $I$ after the value of $J$ is known.
Then, the zero value $S(I|J)=0$ means that $J$ completely predict $I$. 

\subsubsection{Mutual information.}
This can be expressed in various ways. 
It is likely that the name ``mutual information'' was given.
\begin{align*}
  S(I : J) &= S(I) + S(J) - S(I, J) \\
         S(I : J)&= S(I) - S(I|J) = S(J) - S(J|I)
\end{align*}
For example, it represents the reduction of uncertainty about $I$ due to knowing $J$, that is, the increase in knowledge about $I$. 

Expanding explicitly, its relation to the relative entropy becomes visible
\begin{align*}
  S(I : J) &= S(I) + S(J) - S(I, J) \\
&= - \sum_{i=1} p_i \log_2 p_i - \sum_{j=1} q_j \log_2 q_j + \sum_{i,j=1} p(i,j) \log_2 p(i,j) \\
&= - \sum_{ij} p(i,j) \log_2 p_i - \sum_{ij} p(i,j) \log_2 q_j + \sum_{i,j=1} p(i,j) \log_2 p(i,j) \\
S(I : J) 
&= \sum_{i,j=1} p(i,j) \log_2 \frac{p(i,j)}{p_i q_j}
= S(p(i,j)||p_iq_j)
\end{align*}
since $p_i = \sum_j p(i,j)$ and $q_j = \sum_i p(i,j)$.
One sees that $S(I:J)$ represent relative entropy expressing the deviation of the joint probability distribution $\{p(i,j)\}$ from the product distribution of two independent distributions $\{p_i\}, \{q_j\}$.

\subsubsection{Entropy property.}
First is the symmetry and asymmetry
\begin{equation*}
  S(X, Y) = S(Y, X), \quad S(X : Y) = S(Y : X) \qquad S(X|Y) \neq S(Y|X)
\end{equation*}
Then the inequality 
\begin{equation*}
  S(X) \leq S(X,Y)
\end{equation*}
When the goal is \( Y \), it depends only on the relationship with \( X \). 
The interpretation of this relationship can be seen in the mutual information and the proof of its validity.

\subsubsection{Examlpe: Average entropy.}
Consider a coin which is said to come up heads 0 and tails 1 with equal probability, that is $q_0 = q_1 = 1/2$; but in reality is finely adjusted so that it always comes up heads 0, that is $p_0 = 1$ and $ p_1 = 0$.
The true probability distribution is \( \{p_i\} \), and the assumed distribution is \( \{q_i\} \). 
Then, and the relative entropy becomes
\begin{equation*}
  S(p||q) = - p_0 \log_2 q_0 - p_1 \log_2 q_1 + p_0 \log_2 p_0 + p_1 \log_2 p_1 = 1
\end{equation*}

On the other hand, if the coin is said to always come up heads, but in reality is a fair coin that comes up heads and tails with equal probability, then \( p_0 = p_1 = 1/2, q_0 = 1, q_1 = 0 \), and the relative entropy is
\begin{equation*}
  S(p\|q) = - p_0 \log_2 q_0 - p_1 \log_2 q_1 + p_0 \log_2 p_0 + p_1 \log_2 p_1 = \infty
\end{equation*}
In the first scenario, all events that occur under $p$ are assigned nonzero probability by $q$. 
Hence, the divergence is finite.
In the second scenario, the assumed model $q$ assigns zero probability to an event that actually occurs under $p$. 
This leads to a logarithmic divergence, which corresponds to the statement that the prior assumption is completely false, it is falsified immediately by observation.

\subsubsection{Example: Mutual Information.}
Suppose the joint probability distribution $\{p_{ij}\}$ is expressed by the following matrix:
\begin{align*}
  p(i,j) = [p_{ij}] =
\begin{bmatrix}
p_{11} & p_{12} \\
p_{21} & p_{22}
\end{bmatrix}
=
\begin{bmatrix}
1/2 & 0 \\
0 & 1/2
\end{bmatrix}
\end{align*}
In this case, the probability distributions are
\begin{equation*}
  p_i = \sum_j p_{ij} =
\begin{bmatrix}
1/2 \\
1/2
\end{bmatrix}\qquad
q_j = \sum_i p_{ij} =
\begin{bmatrix}
1/2 & 1/2
\end{bmatrix}
\end{equation*}
Therefore, since $S(I) = S(J) $, the mutual information is
\begin{equation*}
  S(I : J) = 
  -2\left(\frac{1}{2}\log_2\frac{1}{2}+\frac{1}{2}\log_2\frac{1}{2}\right)
  +\left(\frac{1}{2}\log_2\frac{1}{2}+\frac{1}{2}\log_2\frac{1}{2}\right)
  =1
\end{equation*}
Since the only possible outcomes are $(I,J)=\{(1,1);(2,2)\}$, if 1 bit of $J$, the second bit, information is known, then $I$, the first bit, information can be completely known as well.

\subsubsection{Example: Mutual Information (Complete independent correlation).}
Consider the following joint probability distribution \( p_{ij} \):
\begin{equation*}
  [p_{ij}] =
\begin{bmatrix}
1/4 & 1/4 \\
1/4 & 1/4
\end{bmatrix}
\end{equation*}
In this case, the probability distributions are
\begin{equation*}
  [p_i] = \sum_j p_{ij} =
\begin{bmatrix}
1/2 \\
1/2
\end{bmatrix}, \qquad
[q_j] = \sum_i p_{ij} =
\begin{bmatrix}
1/2 & 1/2
\end{bmatrix}
\end{equation*}
Like before \( S(I) = S(J)  \), so the mutual information is
\begin{equation*}
  S(I : J) =2 -4\left(\frac{1}{4}\log_2\frac{1}{4}\right)= 0
\end{equation*}
This indicates that even if one knows 1 bit of information about \( J \), no information about \( I \) can be obtained.

\subsubsection{Example: ABCD model}
Examples 1 and 2 make the meaning of mutual information clear. 
Let the alphabet be \( A = [00], B = [10], C = [01], D = [11] \), and assign occurrence probabilities \( 1/2, 1/4, 1/8, 1/8 \), respectively. 
Then the occurrence probabilities of \( A, B, C, D \) can be written in matrix form as follows.
\begin{equation*}
  [p_{ij}] =
\begin{bmatrix}
p_{00} & p_{01} \\
p_{10} & p_{11}
\end{bmatrix}
=
\begin{bmatrix}
1 / 2 & 1 / 8 \\
1 / 4 & 1 / 8
\end{bmatrix}
\end{equation*}
Then
\begin{equation*}
  [p_i] = \sum_j p_{ij} =
\begin{bmatrix}
5 / 8 \\
3 / 8
\end{bmatrix}\qquad
[q_j] = \sum_i p_{ij} =
\begin{bmatrix}
3 / 4 & 1 / 4
\end{bmatrix}
\end{equation*}
We obtain the mutual information
\begin{equation*}
  S(I : J) = S(I) + S(J) - S(I,J) = 0.015
\end{equation*}
This shows that when \( A,B,C,D \) are encoded as \( [00],[10],[01],[11] \), and only the second entry is measured, the extent to which the first entry can be inferred is approximately \( 0.015 \), which is small but nonzero.

On the other hand, if the occurrence probabilities of \( A,B,C,D \) are changed to
\begin{equation*}
  [p_{ij}] =
\begin{bmatrix}
1 / 3 & 1 / 3 \\
1/3& 0
\end{bmatrix}
\end{equation*}
then the mutual information becomes
\begin{equation*}
  S(I : J) = 0.25
\end{equation*}
which is significantly larger. 
In this case, the lower-right entry has zero probability, which contributes strongly to the correlation.

\subsubsection{Key-Search Paradox}
It is assumed that a key exists in a pocket with a probability of $1/2$, and within one of 16 drawers with a probability of $1/2$. 
Assuming all 16 drawers are equivalent, the uncertainty before searching for the key, expressed as Shannon entropy, is derived as
\begin{equation*}
  S(\text{before}) = -\frac{1}{2} \log_2 \frac{1}{2} + 16 \cdot \left( -\frac{1}{32} \log_2 \frac{1}{32} \right) = 3
\end{equation*}
Suppose the pocket is searched, and it is determined that the key is not there. 
In that case, the key must be in one of the drawers, so the Shannon information content is:
\begin{equation*}
  S(\text{after, drawers}) = 16 \cdot \left( -\frac{1}{16} \log_2 \frac{1}{16} \right) = 4
\end{equation*}
Thus, the uncertainty increases compared to before searching the pocket. 
This implies that uncertainty increased because an experiment was conducted, which is perceived as strange. 
This is the reason it is called a paradox.

The problem occurs due to different probability spaces are being mixed. 
This then is resolved by calculating the overall average uncertainty, taking into account the change caused by searching the pocket
\begin{equation*}
  S(\text{after}) = \frac{1}{2} S(\text{after, pocket}) + \frac{1}{2} S(\text{after, drawers})
  = \frac{1}{2} \times 0 + \frac{1}{2} \times 4 = 2
\end{equation*}

When memory is introduced into the information manipulation of the key-searching problem, another perspective becomes apparent. 
First, the search for the key is performed; it is recorded as $P0$ if it is found in the pocket, and as $P1$ if not. 
Based on that record, if $P0$, no further action is taken; if $P1$, the drawers are searched one by one.

If this key-searching process is performed numerous times, the Shannon information content in the memory of the pocket search results becomes $1$. 
Similarly, the Shannon information content in the memory of the drawer search results becomes $4$. 
Finally, since the information entropy of the system becomes zero once the key is found, the sum of the Shannon information content held by the system and the memory is $3$. This is equal to the Shannon information content of $3$ that the system possessed prior to measurement.
Taking the overall average, the Shannon information content of the memory is $1 + \frac{1}{2} \times 0 + \frac{1}{2} \times 4 = 3$.

In this manner, acquiring "knowledge" is the reduction of "uncertainty." 
This "uncertainty" can be quantified using the concept of probability as Shannon information content. 
Ultimately, acquired "knowledge" can be expressed in terms of Shannon information content.
\end{document}